% Source: https://artofproblemsolving.com/community/c3h1570680_3d_geometry
% Stable identity: aops-topic-1570680-q01
The line through $P$ in the direction $(2,1,-2)$ consists of the points
\[
(3+2t,\ -4+t,\ 5-2t).
\]
It meets the plane when
\[
2(3+2t)+5(-4+t)+6(5-2t)=19\qquad\Longrightarrow\qquad 16-3t=19\qquad\Longrightarrow\qquad t=-1,
\]
at the point $Q(1,-5,7)$. Check: $2\cdot1+5\cdot(-5)+6\cdot7=19$.

The required distance is
\[
PQ=\sqrt{(3-1)^2+(-4+5)^2+(5-7)^2}=\sqrt{4+1+4}=\boxed{3}.
\]
(Equivalently, $PQ=|t|\cdot\left|(2,1,-2)\right|=1\cdot3$.)

\begin{asy}
usepackage("amsmath");
size(6cm);
defaultpen(fontsize(10pt));
pair p0=(0,0),p1=(7,0),p2=(9.5,3),p3=(2.5,3);
fill(p0--p1--p2--p3--cycle,lightgray+opacity(0.6));
draw(p0--p1--p2--p3--cycle,linewidth(0.8));
pair Q=(4.5,1.4),P=(6.2,5.6);
draw(P--Q,red+linewidth(1),Arrow(TeXHead));
dot(P);dot(Q);
label("$P(3,-4,5)$",P,NE);label("$Q(1,-5,7)$",Q,S);
label("$3$",(P+Q)/2,W,red);
label("$2x+5y+6z=19$",(8.1,0.3),E);
label("\mbox{}",(min(currentpicture,user=true)+max(currentpicture,user=true))/2,black); // keeps the text after the figure black when LaTeX typesets the labels inline
\end{asy}

(Schematic sketch, not to scale: the segment $PQ$ in the given direction ends on the plane.)
